\PassOptionsToPackage{usenames,dvipsnames,table}{xcolor}
\documentclass[12pt,aps,preprint,nofootinbib,noshowkeys,noshowpacs,longbibliography,superscriptaddress,tightenlines]{revtex4-2}
\usepackage{amsmath,amsfonts,amsthm,amssymb}
%\usepackage{showkeys}
\usepackage{graphics,graphicx}
\usepackage[inline]{enumitem}
\usepackage{mathrsfs}
\usepackage{bm}
\usepackage{mathtools}
\usepackage{bbold}
\usepackage{xcolor}
\definecolor{darkblue}{RGB}{0,0,196}
\definecolor{darkgreen}{RGB}{0,120,0}
\usepackage{tikz}
%\usepackage{cancel}
% \usepackage{caption}
%\usepackage{bbold}
%\usepackage{subfigure}
%\usepackage{stix}
%\usepackage{multirow}
%\usepackage{longtable}
%\usepackage{color}
%\usepackage[normalem]{ulem}
%\usepackage{hyperref}
%\usepackage{authblk}
%\usepackage{bigints}
%\usepackage{xparse}
%\usepackage{physics}
%\usepackage{verbatim}
%\usepackage{minibox}
%\usepackage{comment}
%\usepackage{appendix}
%%%\usepackage[inline]{showlabels}
%\usepackage{scalerel}
%\usepackage{marginnote}
%\usepackage{graphicx}
%\usepackage[nice]{nicefrac}
%\usepackage{soul}
\newcommand{\mathcolorbox}[2]{\colorbox{#1}{$\displaystyle #2$}}
\newcommand{\hlfancy}[2]{\sethlcolor{#1}\hl{#2}}
% -- highlighting
\usepackage{stackengine}
\renewcommand\useanchorwidth{T}
\newcommand\oU[1]{\ensurestackMath{\stackon[1pt]{#1}{\mkern2mu\bullet}}}
\newcommand\oX[1]{\ensurestackMath{\stackon[1pt]{#1}{\mkern2mu\star}}}
\newcommand\oY[1]{\ensurestackMath{\stackon[1pt]{#1}{\mkern2mu\smwhitestar}}}
\newcommand\oZ[1]{\ensurestackMath{\stackon[1pt]{#1}{\mkern2mu\scaleto{\circ}{3pt}}}}
\renewcommand\S{\mathcal S}
\newcommand\Sstress{{\mathcal S}}
\renewcommand\P{\mathcal P}
\newcommand\N{\mathcal N}
\newcommand\uS{{\underline{S}} }
\def\half{\tfrac{1}{2}}
\def\quarter{\tfrac{1}{4}}
\def\LF{{\scriptscriptstyle LF}}
\def\Pbb{{\mathbb P}}
\def\Ebb{{\mathbb E}}
\def\Mbb{{\mathbb M}}
\def\Sbb{{\mathbb S}}
\def\Vbb{{\mathbb V}}
\def\bbeta{{\sf b}}
\def\Teff{{\sf T}}
\newcommand{\ubeta}{{\underline{\smash{\beta}}}}
\newcommand{\uf}{{\underline{{f}}}}
\newcommand{\uu}{{\underline{u}}}
\def\pp{{\rm p}}
\def\three{{}^{(3)}}
\def\F{{\mathcal F}}
\def\K{{\mathcal K}}
\def\x{{\bm x}}
\def\p{{\bm p}}
\def\edense{{\mathcal{E}}}
\def\pdense{M}
\def\dpdense{\delta\kern-0.08em M}
\def\dbeta{\delta\kern-0.08em \beta}
\def\CChi{ {\mathcal \chi}}
\def\HP{\hphantom{\alpha}} % horizontal
\def\llangle{\left\langle}
\def\rrangle{\right\rangle}
\newcommand{\nn}{\nonumber}
\newcommand{\rl}{\mathrm{rel}}
\newcommand{\tbj}{\tau}
\newcommand{\vecxp}{\vec{x}_\perp}
\def\Pr{ {\rm p} }
\def\edense{ {\rm e} }
\def\muh{\hat\mu}
\def\vh{\hat v}
\def\noisekernel{{\kappa}}
\def\Eq#1{eq.~(\ref{#1})}
\def\Eqs#1{eqs.~(\ref{#1})}
\def\eq#1{(\ref{#1})}
\def\app#1{App.~\ref{#1}}
\def\Fig#1{Fig.~\ref{#1}}
\def\Figs#1{Figs.~\ref{#1}}
\def\Sect#1{Sect.~\ref{#1}}
\def\Ref#1{Ref.~\cite{#1}}
\def\betaz{{(\beta_0)}}


% equation environments  
\def\P{\mathcal P}
\def\beq{\begin{equation}}
\def\eeq{\end{equation}}
\def\st{\begin{equation}}
\def\stp{\end{equation}}
\def\ba{\begin{eqnarray}}
\def\ea{\end{eqnarray}}   
\def\DD{{\mathcal L}}
\def\vDD{{\mathcal L}}
\newcommand\vbeta{{\vec{\beta}}}
\newcommand\vcovD{D}
\newcommand\Chi{\mathcal X}
\newcommand\covD{\nabla}
\newcommand{\ve}[2]{e_{\;\,#1}^{#2}}
\newcommand{\de}[2]{e_{\;\,#2}^{#1}}
\newcommand{\vp}{{\vec{p}\,}}
\newcommand{\petsc}{{\tt PETsc}}

\def\E{{\mathcal E}}
\def\P{{\mathcal P}}
\def\M{{\mathcal M}}
\def\m{{\mathfrak m}}
\def\dd{{\rm d}}
\def\lps{{N}}
\def\sp{\phantom{j}}
\def\spm{\phantom{\mu}}
\def\sft{\vec{N}}
\def\tideal{{\mathcal T}}
\def\tvisc{\Pi}
\def\btvisc{\bar{\Pi}}
\newcommand\ph{\phantom{\nu}}

\newcommand\V{V_0}
\newcommand\dtau{{{\Delta}\tau}}
\newcommand\dt{{{\Delta}t}}

%\usepackage{showkeys}
\def\bb1{{(1)}}
\def\bbt{{(2)}}
\def\bbth{{(3)}}

\newcommand{\AM}[1]{{\color{green}AM: #1}}

\newcommand{\cmtJF}[1]{{\color{blue}JF: #1}}
\newcommand{\rs}[1]{{\color{red}RS: #1}}
\newcommand{\cmtMS}[1]{{\color{cyan}MS: #1}}

\newcommand{\gammathree}{{}^{(3)}\Gamma}
\usepackage[colorlinks=true,linktocpage=true,linkcolor=darkblue,citecolor=red,urlcolor=darkblue]{hyperref}
\usepackage{listings}
\definecolor{jlkw}{RGB}{0,0,180}
\definecolor{jlcmt}{RGB}{0,128,0}
\definecolor{jlstr}{RGB}{160,60,0}
\definecolor{jlfn}{RGB}{128,0,128}
\lstdefinelanguage{Julia}{%
  morekeywords=[1]{abstract,break,catch,const,continue,do,else,elseif,%
    end,export,false,for,function,immutable,import,if,in,%
    macro,module,mutable,return,struct,true,try,using,while,with},%
  morekeywords=[2]{zero,dot,minimum_image,wendland_c2_grad,%
    kappa_tensor,neighbors,SMatrix,rebuild_cell_list},%
  sensitive=true,%
  alsoletter={!},%
  morecomment=[l]{\#},%
  morecomment=[n]{\#=}{=\#},%
  morestring=[b]{"},%
}[keywords,comments,strings]

\lstset{
  language=Julia,
  basicstyle=\ttfamily\small,
  columns=flexible,
  keepspaces=true,
  showstringspaces=false,
  keywordstyle=[1]\color{jlkw}\bfseries,
  keywordstyle=[2]\color{jlfn},
  commentstyle=\color{jlcmt}\itshape,
  stringstyle=\color{jlstr},
  literate=
    {ε}{$\varepsilon$}{1}
    {κ}{$\kappa$}{1}
    {Δ}{$\Delta$}{1}
    {β}{$\bbeta$}{1}
    {>=}{$\geq$}{2}
    {!=}{$\neq$}{2}
}



\renewcommand{\baselinestretch}{1.05}

\begin{document}

\section{Field and particle notation}

The continuous entropy density field is denoted $S(x)$; the relativistic entropy density is
\st
s_{\rm rel}(x) \equiv \frac{S(x)}{\gamma(x)} \, .
\stp
Each SPH particle $a$ carries a conserved scalar $s_a$ (entropy per particle).
After smoothing, the SPH estimate of the entropy density at the particle site is
\st
S(x_a) = \sum_b s_b\, W(x_a, x_b) \, .
\stp
The per-particle entropy satisfies
\st
s_a \approx S(x_a)\,\mathrm{vol}_a \, ,
\stp
where $\mathrm{vol}_a = L^2/N$ is the area per particle.  This relation is
\emph{approximate}: $S(x_a)$ is itself a kernel average, so there is an $\mathcal{O}(h^2)$
smoothing error, and for a non-uniform lattice the exact equality never holds.
The setup functions \texttt{set\_ic!} and \texttt{setup\_ic} in \texttt{shearwave.jl}
initialise particles via $s_a = S_{\rm func}(x_a)\,\mathrm{vol}_a$, where
$S_{\rm func}(x,y)$ is the user-supplied continuous entropy density.

\section{Algorithm}
We define the sph kernel $W(x_a, x_b) = F(\half (x_a  - x_b)^2 )$ and the derivative
\st
D_k(x_a, x_b) = \frac{\partial W(x_a, x_b)}{\partial x_a^k} 
\stp

\begin{enumerate}
    \item Given $m_a^i$ and  $s_a$  compute $M^i(x_a), S(x_a)$. 
	\begin{align}
	    M^i(x_a) =& \sum_b m_b^i W(x_a, x_b) \\
	    S(x_a) =& \sum_b s_b W(x_a, x_b)
	\end{align}

    \item Use the equation of state to find  the thermodynamic conjugates,  $(v/S)_i(x_a)$ and  $-(\F/S)(x_a)$,  from the extensive variables $M^i(x_a), S(x_a)$.  
      Also compute the potential $\E(x_a)$ and (optionally) its legendre transform $T(x_a)$. See  \Sect{sec:eos} for details.

\item Compute a number of quantities
    \begin{align}
	\dot x_{i, a} \equiv& \sum_b s_b (v/S)_i (x_b) W(x_a, x_b) \\
	\Teff_a =& \E(x_a)  - \sum_b s_b (\F/S)(x_b) W(x_a, x_b)  \\
	\bbeta_{i, a} \equiv&  \frac{\dot x_{i, a}}{\Teff_a}
    \end{align}
These are related to the derivatives of the Hamiltonian (see below)
\item  
The Hamiltonian is 
\st
{\mathbb H}(s_a, x_a, m_a)  = \sum_a s_a \E(M(x_a), S(x_a))
\stp
and we have (as advertised)
\st
\Delta {\mathbb H} (s_a, x_a, m_a)  =  \sum_a \Teff_a  \Delta s_a  + \dot x_{a} \cdot  \Delta m_a 
\stp

The ideal equations are Hamilton's. Spefically we do the following:
  Update the positions
    \st
    \Delta x^i_a = \Delta t \, \delta^{ik }\dot x_{k, a}
    \stp
    and then update the momentum due to the ideal terms 
    \begin{multline}
\Delta m_{a}^i =  -\Delta t \sum_{b} \big[ (v/S)_j(x_a) s_a m_b^j + (v/S)_j(x_b) s_b m_a^j  \\
    -  (\F/S)(x_a)s_a s_b - (\F/S)(x_b) s_b s_a \big] \delta^{ik }  D_k (x_a, x_b) \, . 
\end{multline}
\item Update the momentum and entropy due to viscous forces. Update the momentum by viscous forces 
    \begin{align}
      \Delta m^i_{{\rm visc}, a} =& \Delta t  \, f^i_{\rm visc}(x_a) \\
      \Delta s_{a} =& -\bbeta_a  \cdot \Delta m_{{\rm visc}, a}
   \end{align}
   The viscous force is assembled in two passes --- the two-level scheme of
   \eqref{eq:strain}, \eqref{eq:K_twolevel} and \app{app:dissipative_loop}.
   First compute the SPH strain at every particle,
   \st
   \varepsilon_{ji}(x_a) = \frac{1}{S(x_a)} \sum_b s_b\,
      \big(\bbeta_{i,b}-\bbeta_{i,a}\big)\, D_j(x_a,x_b) \,,
   \stp
   then assemble the force from the stored strains,
   \begin{gather}
   f^i_{\rm visc}(x_a) = \sum_{b} F^i_{\rm visc}(x_a, x_b) = - \frac{\partial K}{\partial \bbeta_{i, a}} \, ,  \\
   F_{\rm visc}^i(x_a, x_b) =
   s_a s_b\, D_j(x_a, x_b) \left(
      \frac{T_{\rm rel} \kappa^{ijlm}}{S^2}\bigg|_a \varepsilon_{lm}(x_a)
      + \frac{T_{\rm rel} \kappa^{ijlm}}{S^2}\bigg|_b \varepsilon_{lm}(x_b)
      \right) \,.
   \end{gather}
   The total rate of entropy production is
   \st
   \frac{d\Sbb}{dt} \;=\; 2K \;=\; \sum_a \frac{s_a}{S(x_a)}\;
      \varepsilon_{ji}(x_a)\, T_{\rm rel}\kappa^{ijlm}\big|_a\,
      \varepsilon_{lm}(x_a) \;\ge\; 0\,,
   \stp
   with $K$ the dissipative function \eqref{eq:K_twolevel}.


    \end{enumerate}
\section{\bf Stochastic updates} 

First we  generate the noise for each pair 
\begin{align}
    \llangle \xi^{ij} \xi^{lm}  \rrangle = 2 T_{\rm rel} \kappa^{ijlm}(v) \delta(t-t') \delta^2(\mathbf{r}-\mathbf{r}')
\end{align}

\section{Hamiltonian evolution}

\noindent{\bf Starting point}
\st
\label{eq:Hamiltonian}
\boxed{
{\mathbb H}(s_a, x_a, m_a)  = \sum_a s_a \E(\M(x_a), V_s(x_a))
}
\stp
The variables in the Hamiltonian are summarized below. \\

\noindent {\bf Thermo: }
We have
\st
d\Ebb = Td{\Sbb} + v d\Mbb - p d\Vbb
\stp
We work with a volume where $\Sbb = \Sbb_0$ is constant. Then $S= \Sbb_0/\Vbb$.  Divinding by 
$\Sbb_0$ and definining $d\Vbb/\Vbb = -dS/S(x)$ one has 
\st
 d \E =  v d\M   - p d(1/S)
\stp
where $\E= \Ebb/\Sbb$ and $\M = \Mbb/\Sbb$.   We 
also note 
\st
d \E =  \frac{v}{S} dM  - \frac{\F}{S}\, dS
\stp
where the free energy per entropy is $\F$:
\st
\mathcal F \equiv \E - T =  v M/S - p/S   \qquad \F/S \equiv v M/S^2 - p/S^2
\stp

We also note that $T\equiv-1/\beta_0(x)$ is the related to the ``usual" $T_{\rm rel}$:
\st
T \equiv \frac{1}{\beta^0} \equiv \frac{T_{\rm rel}}{\gamma}.
\stp
% From now on we define $V_s$ as the specific volume $1/s$
% \st
%  d \E =  v d\M   - p dV_s
% \stp
%
% And a convenent legendre transform
% \st
% d (\E + p V_s) =  v d\M +  V_s dp
% \stp
% and 
% \st
% d (\E + p V_s - v \M) = dT  =  -\M  dv  + \frac{1}{s} dp
% \stp

\noindent{{\bf Variables: }}
The variables are
\begin{align}
    S(x) = \frac{1}{V(x)}  \equiv& \sum_a s_a W(x, x_a) \\
    M(x) \equiv & \sum_a  m_a W(x, x_a)  \\
    \M (x)  \equiv& \frac{1}{S(x)} \sum_{a} m_a W(x, x_a)  \\
    m_a \equiv& \m_a s_a
\end{align}
The conjugate variables are functions of these, i.e. 
\st
   v(x_a) = v(\M(x), 1/S(x)) \neq v(\M_a, 1/s_a)
\stp
% \st
%  \partial \M(x) /\partial M_a =  s_a W(x, x_a)/S(x)
% \stp
% \begin{align}
%     \partial \M(x) /\partial x_a =&  -s_a \M_a W'(x, x_a)/S(x) -   \frac{M(x)}{S^2(x)} (-s_a  W'(x, x_a)) 
% = -s_{a} W'(x, x_a)/S(x) \left( \M_a  - \M(x)  \right)
% \end{align}
\noindent {\bf EOM}: 
We use 
\begin{align}
    \Aboxed{
    \frac{dx_a}{dt}  =&  \frac{\partial {\mathbb H}}{\partial m_a} }  \\
    % = \frac{1}{s_a}  \frac{\partial {\mathbb H}}{ \partial \M_a}  \\
    % =&  \sum_b s_b v(x_b) \frac{\partial \M(x_b)}{\partial \M_a}   \\
    =& \sum_b s_b v(x_b) W(x_b, x_a) /S(x_b) \\
   % =&  \frac{1}{V(x_a)} \underline{V v}(x_a)  \\
    \equiv& S(x_a) \underline{v/S}(x_a)
\end{align}
where we use the underline to denote an average:
\st
\underline{Q}(x) \equiv \frac{1}{S(x) } \, \sum_b\,  Q_b \,  s_b \, W(x_b, x)  
%\equiv  \sum_{b} Q_b \lambda_b(x)    \qquad  \lambda_b(x) \equiv  s_b W(x_b, x)/S(x)
\stp
Thus
\st
\dot x_b - v(x_b)  \equiv  S(x_b) \underline{v/S}(x_b) - v(x_b)
\stp
And
\st
\M(x) = \underline \m(x) = \underline{m/s}(x)
\stp


We also have
\st
\boxed{
    \frac{dm_a}{dt} =  -\frac{\partial {\mathbb H}}{\partial x_a}  +  f^{i}_{\rm visc}(x_a) 
}
\stp
Computing the derivatives we have 
\begin{multline}
    \frac{dm_{a}^i}{dt}  =  -\sum_{b} s_a s_b \left[ (v/S)^j(x_a) \Delta \M_b^j(x_a) + (v/S)^j(x_b) \Delta \M_a^j (x_b)  \right]  \delta^{ik } D_k (x_a, x_b)
    \\
    - \sum_{b} s_a s_b \left[ (p/S^2)(x_a) + (p/S^2)(x_b) \right] \delta^{ik }  D_k (x_a, x_b)
    + f_{\rm visc}^i(x_a)
\end{multline}
Here 
\st
\Delta \M_b(x_a) = \m_b - \M(x_a)
\stp

We also defined the derivative operator
\st
 D_{k}(x_a,x_b) \equiv \frac{\partial W(x_a , x_b) }{\partial x_a^k }
\stp

In retrospect it seems natural to write the force term as 
\begin{multline}
\frac{dm_{a}^i}{dt}  =  -\sum_{b} \left[ (v/S)^j(x_a) s_a m_b  + (v/S)^j(x_b) s_b m_a  \right]  \delta^{ik } D_k (x_a, x_b)
    \\
    + \sum_{b} \left[ (\F/S)(x_a)s_a s_b + (\F/S)(x_b) s_b s_a \right] \delta^{ik }  D_k (x_a, x_b)
    + f_{\rm visc}^i(x_a)
\end{multline}

\noindent {\bf Proposed viscous forces:} 
When viscous forces are included we include the last line which 
we write just like the pressure.
\begin{align}
    \Aboxed{
    f_{\rm visc}^i(x_a) \equiv& -\sum_{b} s_a s_b \left( (\Pi^{ik}/S^2)(x_a) + (\Pi^{ik}/S^2)(x_b) \right)  D_k (x_a, x_b)
} \\
    \equiv& \sum_{b} F_{\rm visc}^i(x_a, x_b)
\end{align}
Then net momentum transfer is zero 
\st
\sum f_{\rm visc}(x_a)= 0 
\stp

\noindent {\bf Change in entropy: }
Let's see how the energy changes under viscous evolution. 
We use the change in the entropy written as
\begin{align}
    \Delta S_{\rm visc} (x_a) =& \sum_{b} \Delta s_b   W(x_a, x_b )  % = S(x_a) \underline{ (\Delta s/s) }(x_a ) 
\end{align}
Starting from \Eq{eq:Hamiltonian}  we find
\st
\Delta {\mathbb H} = \sum_{a}\left[ \Delta s_a \, \E(\M, V^s)   - s_a \, (\F/S)_a  \Delta S_{\rm visc}(x_a)  \right]   + \sum_{a} \frac{\partial {\mathbb H}}{\partial m_a}  \Delta m_{\rm visc, a}
\stp
% We have
% \st
% {\mathcal F}(x)\equiv \E(x) - T(x)  = v M/S  -p/S = \mbox{Free energy per entropy} 
% \stp

% Then lets look at the energy density in each cell
% \st
% E(x_a) = S(x_a) \E( M(x_a)/S(x_a), 1/S(x_a)) 
% \stp
Then noting
\st
\sum_{a} s_a \, Q(x_a) \, \Delta S(x_a) =  \sum_{ab}  \Delta s_b  \,  W(x_b, x_a)  \, s_a Q(x_a) = \sum_b \Delta s_b \, S(x_b) \,  \underline{Q}(x_b)
\stp
the free energy term can be written
\st
\sum_a  s_a (\F/S)_a \, \Delta S_{\rm visc}(x_a)   = \sum_b \Delta s_b \, S(x_b)  \, \underline{{\mathcal F}/S}(x_b)
\stp
So we have 
\st
\Delta {\mathbb H} = \sum_a \Delta s_a  \, \Teff_a +  \frac{dx_a}{dt}  \Delta m_{\rm visc, a} = 0
\stp
where here: 
\st
\boxed{
\Teff_a = T(x_a)  + \left({\F}(x_a) - S(x_a) \underline{\F/S}(x_a) \right)
}
\stp
So to have energy conservation we have 
\st
\boxed{
    \Delta s_a =  -\frac{1}{\Teff_a} \frac{dx_a}{dt} \cdot \Delta t \, f_a 
}
\stp
where we said $\Delta m_{\rm vis, a} = \Delta t f_a$.  Then 
\st
\boxed{
\Delta \Sbb = \sum_a \Delta s_a = - \sum_{ab} \frac{\dot x_a}{\Teff_a}  \cdot \Delta t\, F_{ab}\ =  - \sum_{ab} \left( 
    \frac{\dot x_a}{\Teff_a}   
    -
\frac{\dot x_b}{\Teff_b}   \right) \cdot \Delta t \, F_{\rm visc}(x_a, x_b) 
}
\stp

 
\noindent {\bf Entropy production,  deterministic hydro, and the dissipative function } \\

We start by postulating the dissipative function
Define the SPH strain at particle $a$ as
\begin{equation}
\label{eq:strain}
    \varepsilon_{ji}(x_a) \;\equiv\;
    \frac{1}{S(x_a)}\sum_b s_b\,
    \bigl(\bbeta_{i,b} - \bbeta_{i,a}\bigr)\,D_j(x_a,x_b) \,,
\end{equation}
which is an antisymmetric SPH gradient estimator for $\partial_j\beta_i$ at
particle $a$.  The discrete dissipative function is then
\begin{equation}
\boxed{
\label{eq:K_twolevel}
    K \;=\; \frac{1}{2} \sum_a \frac{s_a}{S(x_a)}\;
    \varepsilon_{ji}(x_a)\,
    T_{\rm rel}\kappa^{ijlm}\big|_a\;
    \varepsilon_{lm}(x_a) \,.
}
\end{equation}
Here the two factors of $D_j$ live in \emph{separate} neighbor sums (one for
each index of $\varepsilon$), so there is no overlap.  The volume element
$s_a/S(x_a)$ is the correct SPH approximation to $d^2x$.
We have defined
\st
\bbeta_i(a) \equiv \frac{\dot x_{a, i} }{\Teff_a}
\stp
As we will see the rate of entropy production is 
\st
 \dot \Sbb = 2 K 
\stp

We take 
\st
\boxed{
f_{\rm visc}^i(x_a)  = -\frac{\partial K}{\partial \bbeta_{i, a }  } = \sum_b F^i_{\rm visc}(x_a, x_b) }
\stp
where the pair of forces is
\st
\boxed{
      F^i_{\rm visc}(x_a, x_b) = s_a s_b \,D_j(x_a, x_b)\,
     \left( \left. \frac{T_{\rm rel}\kappa^{ijlm}}{S^2}\right|_a\;\varepsilon_{lm}(x_a)
    \;+\;
    \left. \frac{T_{\rm rel}\kappa^{ijlm}}{S^2}\right|_b\;\varepsilon_{lm}(x_b) \right)
}
\stp
We note that
\st
 F^i_{\rm visc}(x_a, x_b) = - F^i_{\rm visc}(x_b, x_a)
\stp
due to the antisymmetry of $D_i(a,b)$. 
This fits into the general mould where momentum is conserved since
\st
\sum_{a} f_{\rm visc}^i(x_a) =\sum_{ab} F^i(x_a, x_b)= 0 
\stp

So the previously defined mean strains are now more explicit:
\st
   \left(\frac{\Pi^{ij}}{S^2}(x_a)  
   +
   \frac{\Pi^{ij}}{S^2}(x_b)  
   \right)
     =
     -\left( \left. \frac{T_{\rm rel}\kappa^{ijlm}}{S^2}\right|_a\;\varepsilon_{lm}(x_a)
    \;+\;
    \left. \frac{T_{\rm rel}\kappa^{ijlm}}{S^2}\right|_b\;\varepsilon_{lm}(x_b) \right)
\stp
with 
\begin{align}
    \Aboxed{
    F_{\rm visc}^i(x_a, x_b) \equiv& - s_a s_b \left( (\Pi^{ik}/S^2)(x_a) + (\Pi^{ik}/S^2)(x_b) \right)  D_k (x_a, x_b)
} 
\end{align}




The signs and magnitude are consistent with the  continous expression 
\st
\Pi^{ij} = - T_{\rm rel} \kappa^{ijlm} \partial_l \beta_m
\stp
The net force on a particle  
\st
f^i_{\rm visc}(x_a) = - \sum_b s_a s_b  D_j (x_a ,x_b) \left(\frac{\Pi^{ij}}{S^2} (x_a)  + \frac{\Pi^{ij}}{S^2} (x_b) \right)   
\stp
This 
is compatible with the continous expression 
\st
  \partial_t M^i = - \partial_j \Pi^{ij}
\stp

Finally and most importantly
\st
\boxed{
\dot \Sbb =  \sum_{a} -\bbeta_a \cdot f_{\rm  visc}(x_a) =   \bbeta_a \cdot  \frac{\partial K }{\partial \bbeta_{a} } =\,  2  K \geq 0
}
\stp



% We also note that
% \st
% \frac{\Delta S_{\rm visc}}{S} (x_a) = \sum_{b} (\Delta s_b/s_b)  s_b W(x_a, x_b ) =  \underline{ (\Delta s/s) }(x_a )
% \stp
% Let look at
% \begin{multline}
%     \Delta E(x_a) =   S(x_a) \Big( \frac{\Delta S_{\rm visc}(x_a)}{S(x_a) } \E(x_a)  +  v(x_a) \Delta M_{\rm visc}(x_a)/S(x_a) \\ - v(x_a) \frac{M(x_a)}{S(x_a)}  \frac{\Delta S_{\rm visc}}{S}(x_a)  + \frac{p(x_a)}{S(x_a) } \frac{\Delta S_{\rm visc} }{S(x_a) } \big] 
% \end{multline}
% Using 
% \st
%   \E - v M/S + p/S = T(x)
% \stp
% we find
% \st
% \Delta E(x_a) = T (x_a) S(x_a) \big[  (\frac{\Delta S(x_a)}{S(x_a)}   + \frac{v(x_a)}{T(x_a) }  \frac{\Delta M_{\rm visc}(x_a) }{S(x_a)}
% \stp
% So if we require that $\Delta E(x_a) = 0$,  then  
% \st
% \Delta s_a = \frac{v(x_a) }{T(x_a) } \Delta M_{a, \rm visc} 
% \stp


\subsection{Some notes} 

\subsubsection{Translational Invariance}
We have charges
\st
Q(x) = \sum_a q_a W(x, x_a)
\stp
We have 
\st
W(1, 2) = F(\tfrac{1}{2} (\vec{x}_1 -\vec{x}_2)^2 ) = F(1, 2)
\stp
\begin{align}
    \frac{\partial W(1,2)}{\partial \vec{1}}   =& F'(1, 2) \, (\vec{x}_1 - \vec{x}_2) \equiv W'(1,2 ) \\
    \frac{\partial W(1, 2)}{\partial \vec{2}} =&  -F'(1, 2) (\vec{x}_1 - \vec{x}_2) \equiv  - W'(1,2 )
\end{align}
So
\st
\frac{\partial Q(x_b) }{\partial x_a}  = \left(\sum_{c} q_c W'_{ac}\right)\delta_{ab} - q_a W'_{ba}  \qquad \qquad \mbox{no-$a$-sum! } 
r 
\stp
We note that  translational invariance implies
\st
 \sum_a  \frac{\partial Q(x_b) }{\partial x_a} = 0 
\stp

\subsubsection{\texorpdfstring{Derivatives with respect to $x_a$}{Derivatives with respect to x\_a}}
We have the variables 
\st
 Q^I =  (M, S)
\stp
Then 
\st
\E_I \equiv \frac{\partial \E}{\partial u^I} =  (\frac{v}{s} ,  -\F/S)
\stp
\begin{align}
    \frac{\partial \mathbb H}{\partial x_a } =& \sum_b s_b \frac{\partial \E}{\partial x_a} \\
    =& \sum_{b} s_b \E_I(x_b) \frac{\partial Q^I(x_b)}{\partial x_a}  \\
    =& \sum_b s_b\E_I(x_b) \left[  (\sum_c q_c^I W'_{ac}) \delta_{ba} -  q_a^I W'_{ba}  \right]   \\
    =& \sum_{c} s_a \E_I(x_a) q^I_c W'_{ac}   - \sum_{b} q_a^I W'_{ba} s_b \E_I(x_b)  \\
    =& \sum_{b}  \left[ s_a \E_I(x_a) \; q^I_{b} + s_b \E_I(x_b) \; q_a^I \right] W'_{ab}
\end{align}

\subsubsection{\texorpdfstring{Evaluation of $\partial \mathbb H/\partial x_a$}{Evaluation of dH/dx\_a}}
For the specific case we have 
\begin{multline}
   \frac{\partial \mathbb H}{\partial x_a} = 
   \sum_{b} \left[  s_a\left( (v/S)(x_a) m_b -  (vM/S^2)(x_a)s_b \right) + s_b \left( (v/S)(x_b) m_a -  (vM/S^2)(x_b) s_a  \right) \right]  W'_{ab} \\ + \mbox{pressure}
\end{multline}
Defining 
\st
 \Delta \M_b(x_a) \equiv \m_b -  \M(x_a) 
\stp
We have
\st
   \frac{\partial \mathbb H}{\partial x_a} = 
   \sum_{b}   s_a s_b \left[   \, v/S(x_a) \, \Delta \M_b(x_a) +  \, v/S(x_b)\,  \Delta \M_a(x_b) \right]  W'_{ab} + \mbox{pressure}
\stp
The pressure term is 
\begin{align}
    \mbox{pressure} =& 
    \sum_{b} \left[ s_a (p/S^2)(x_a) s_b + s_b (p/S^2)(x_b) s_a \right] W'_{ab} \\
    =& 
    \sum_{b} s_a s_b \left[  (p/S^2)(x_a) +  (p/S^2)(x_b) \right] W'_{ab}  
\end{align}
So to summarize 
\st
\frac{\partial \mathbb H}{\partial x_a} = 
\sum_{b}   s_a s_b \left[   \, v/S(x_a) \, \Delta \M_b(x_a) +  \, v/S(x_b)\,  \Delta \M_a(x_b)  +  (p/S^2)(x_a) + (p/S^2)(x_b) \right]  W'_{ab} 
\stp

\subsubsection{Solving the equation of state}
\label{sec:eos}



We have the momentum $M^i(x) = (e + p) \gamma^2 v^i$ and the entropy
\st
    S = \frac{(e+p)\,\gamma}{T_{\rm rel}(e)} \, ,
\stp
where $T_{\rm rel}(e)$ is the relativistic temperature determined by the
equation of state.  Given $S_{\rm in} \equiv S(x_a)$ and
$M \equiv |\vec M(x_a)|$, we want to find $e$ and hence the thermodynamic
conjugates $(v/S)$ and $-(\F/S)$:
\st
 \dd \E = (v/S)_i \, \dd M^i  -  (\F/S) \, \dd S \, ,
\stp
where $\E$ is the energy per entropy.
The spatial index is suppressed below, writing $M^i = M n^i$ and
$(v/S)_i = (v/S)\, n_i$ with $n^i$ a unit vector.

\noindent{\bf Step 1: velocity from a trial $e$.}
From $M^i = (e+p)\gamma\,u^i$ with $u^i = \gamma v^i$, taking magnitudes:
\st
    u = \frac{M}{e+p(e)} \, , \qquad
    \gamma = \sqrt{1 + u^2} \, , \qquad
    v = \frac{u}{\gamma} \, .
\stp
This is exact given $e$; no sub-iteration is required.

\noindent{\bf Step 2: Newton iteration on $e$.}
Define the residual
\st
    f(e) = \frac{(e+p(e))\,\gamma(e)}{T_{\rm rel}(e)} - S_{\rm in} \, ,
\stp
where $\gamma(e)$ is evaluated via Step~1.
The Newton derivative follows from the thermodynamic identity
\st
    T_{\rm rel}\frac{de}{dT_{\rm rel}} = \frac{e+p}{c_s^2}
    \quad\Longleftrightarrow\quad
    \frac{dT_{\rm rel}}{de} = \frac{c_s^2\,T_{\rm rel}}{e+p} \, ,
\stp
and the lemma $d[(e+p)\gamma]/de = (1+c_s^2)/\gamma$ proved below:
\begin{align}
    f'(e) &= \frac{1}{T_{\rm rel}}\frac{d[(e+p)\gamma]}{de}
             - \frac{(e+p)\gamma}{T_{\rm rel}^2}\frac{dT_{\rm rel}}{de}
           = \frac{1+c_s^2}{T_{\rm rel}\,\gamma} - \frac{c_s^2\,\gamma}{T_{\rm rel}}
           = \frac{1 - c_s^2 u^2}{T_{\rm rel}\,\gamma} \, .
\end{align}
The Newton update is $e \leftarrow e - f(e)/f'(e)$.

\noindent{\bf Lemma.}  $\dfrac{d}{de}[(e+p)\gamma] = \dfrac{1+c_s^2}{\gamma}$.

Let $w = e+p$ and $u = M/w$.  Then $dw/de = 1+c_s^2$ and
$du/de = -(u/w)(1+c_s^2)$, so
$d\gamma/de = (u/\gamma)\,du/de = -u^2(1+c_s^2)/(\gamma w)$.  Hence
\st
    \frac{d(w\gamma)}{de}
    = (1+c_s^2)\gamma - \frac{u^2(1+c_s^2)}{\gamma}
    = \frac{1+c_s^2}{\gamma}\underbrace{(\gamma^2 - u^2)}_{=\,1}
    = \frac{1+c_s^2}{\gamma} \, .
\stp

\noindent{\bf Post-convergence expressions.}
With $S \equiv S_{\rm in}$:
\begin{align}
    p =& \; p(e) \, , \\
    \frac{v}{S} =& \; \frac{v}{S} \, , \\
    \frac{\F}{S} =& \; \frac{v}{S}\frac{M}{S} - \frac{p}{S^2} \, , \\
    \E =& \; \frac{(e+p)\gamma^2 - p}{S} \, , \\
    T =& \; \frac{e+p}{S} \, .
\end{align}

\section{Dissipative Function}

We use the Hamiltonian formalism of  {\tt DZYALOSHINSKII AND VOLOVICK } section 2:
\begin{align}
    T \left(\partial_t S(x)  +\{S(x), H \} \right) =&     
    T \partial_\ell \left(\frac{\partial R }{ \partial \nabla_\ell T  }   \right)  + T K \\         
    v^k \left(\partial_t M_k(x)  + \{M_k(x), H \}  \right)=&   
v^k \partial_\ell \left( \frac{\partial R }{  \partial \nabla_\ell v^k  }  \right)
\end{align}
Here $R=K=0$ corresponds to ideal hydro.  $K$ is rate of entropy production.   As we will see,  $R = T K/2$, and $R$ is the dissipative function. In the density frame we know that
\st
K = T_{\rm rel} \kappa^{ijlm} \partial_i\beta_j \partial_l \beta_m
\stp
One of the goals of this section is to relate $K$ to the density frame entropy current. 

To see the relation between $R$ and $K$ we have
\begin{align}
    \partial_t E  + \partial_k \tideal^{0k} =& T 
    \partial_\ell \left(\frac{\partial R }{ \partial \nabla_\ell T  }   \right)  
+    v^k \partial_\ell \left(\frac{\partial R }{ \partial_\ell v^k  } \right) + T K \\
    =&  \partial_{\ell} \left( T \frac{\partial R }{ \partial \nabla_\ell T  }  + v^k  \frac{\partial R}{\partial \nabla_\ell v_k }
	\right)   
	-
      (\nabla_\ell T) \left(\frac{\partial R }{ \partial \nabla_\ell T  }   \right)  
      - (\nabla_\ell v^k )  \left(\frac{\partial R }{ \partial \nabla_\ell v^k  } \right) + T K
\end{align}
Demanding that the rhs is zero up to a total divergence, we find
the relation between $R$ and $K$.
In general the dissipative functional involves
\st
2 R  = T K = T T_{\rm rel} \kappa^{ijlm} \partial_{i}\beta_j  \partial_{\ell} \beta_m
\stp
where
\st
\partial_{(i} \beta_{j)} = \frac{1}{T} \left( -\frac{v_{(i}\partial_{j)}T }{T} + \partial_{(i} v_{j)} \right)  
\stp
Then,   the dissipative flux to the energy
\st
\partial_{\ell} \left( T \frac{\partial R }{ \partial \nabla_\ell T  }  + v^k  \frac{\partial R}{\partial \nabla_\ell v_k }
\right) = 0 
\stp
This is the density-frame condition


The  strain according to us  is 
\st
\Pi^{ij} = -T_{\rm rel} \kappa^{ijlm} \partial_{(\ell} \beta_{m)} 
\stp
which does indeed work:
\st
\Pi^{ij}_R = -\frac{ \partial R}{\partial \nabla_\ell  v_k }  = - T_{\rm rel} \kappa^{ijlm} \partial_{(\ell} \beta_{m)} \, .
\stp
Then we have according to us:
\st
\partial_t S + \partial_i (S v^i) = \partial_j (\beta_i \Pi^{ij})   +   \partial\beta \cdot T_{\rm rel} \kappa \cdot \partial\beta 
\stp
We identify the entropy flux from the dissipative function
\st
\frac{\partial R}{\partial( \nabla_j T) } = -v_i \kappa^{ijlm} \partial_l\beta_m = \beta_i \Pi^{ij}
\stp
which is in agreement. 

We see that the entropy conservation law can be written
\st
\partial_t S + \partial_i (S v^i)= \beta_{i} \,  \partial_{j} \Pi^{ij} 
\stp
Dividing by $N$ and using the conservation law
\st
 \partial_t N + \partial_i N v^i  = 0
\stp
\st
D_t (S/N)  = \frac{1}{N(x)} \beta_i \partial_j \Pi^{ij}
\stp
One can take $N=S^{(0)}(t, x)$ the initial entropy.

\subsection{\texorpdfstring{The $\kappa^{ijlm}$ tensor}{The kappa\^{}\{ijlm\} tensor}}
\label{sec:kappa}

We use the $d=3$ formula for $\kappa^{ijlm}$ throughout, even though the SPH
simulation is restricted to 2D motion in the $(x,y)$ plane.  This choice is
physically natural: the fluid is genuinely three-dimensional and the simulation
is a 2D cross-section.  The $d=3$ formula gives the standard coefficient
$\kappa^{xxxx}|_{v=0} = 4\eta/3$ in the non-relativistic limit.

The tensor splits into shear and bulk parts,
\st
T_{\rm rel}\kappa^{ijlm}
   = \kappa_{\rm shear}^{ijlm}(v)\,\frac{T_{\rm rel}\,\eta}{(1-c_s^2 v^2)^2}
   + \kappa_{\rm bulk}^{ijlm}(v)\,\frac{T_{\rm rel}\,\zeta}{(1-c_s^2 v^2)^2}\,,
\stp
where $i,j,l,m\in\{x,y\}$ are purely spatial indices and
$v^2 = v_x^2 + v_y^2$.  Define the spatial projector
\st
L^{ij} = \delta^{ij} - v^i v^j
\stp
and the $d=3$ abbreviations
\begin{align}
\langle PL\rangle &= 1 - c_s^2 v^2 \,,\\
(LPL)^{ij} &= -\frac{c_s^2\,v^i v^j}{\gamma^2} + \frac{L^{ij}}{3} \,,\\
\langle PLPL\rangle &= \frac{2}{3}(c_s^2 v^2)^2 + \frac{1}{3}(1-c_s^2 v^2)^2 \,.
\end{align}

\noindent{\bf Shear part:}
\st
\kappa_{\rm shear}^{ijlm} = 2\langle PL\rangle^2\left[L^{(il}L^{j)m}
  - \frac{L^{ij}(LPL)^{lm}}{\langle PL\rangle}
  - \frac{L^{lm}(LPL)^{ij}}{\langle PL\rangle}
  + \frac{\langle PLPL\rangle}{\langle PL\rangle^2}\,L^{ij}L^{lm}\right]
\stp
where $L^{(il}L^{j)m} = \frac{1}{2}(L^{il}L^{jm}+L^{im}L^{jl})$.

\noindent{\bf Bulk part:}
\st
\kappa_{\rm bulk}^{ijlm} = L^{ij}\,L^{lm}
\stp
For a conformal EOS, Weyl invariance requires $\zeta=0$, so the bulk term vanishes.

\noindent{\bf Checks at $v=0$:}
\begin{align}
T_{\rm rel}\kappa^{xxxx}\big|_{v=0} &= \tfrac{4}{3}\,T_{\rm rel}\eta\,,\\
T_{\rm rel}\kappa^{xxyy}\big|_{v=0} &= -\tfrac{2}{3}\,T_{\rm rel}\eta\,,\\
T_{\rm rel}\kappa^{xyxy}\big|_{v=0} &= T_{\rm rel}\eta\,.
\end{align}

\noindent{\bf 1D check} ($v_y=0$, speed $v=v_x$):
\st
T_{\rm rel}\kappa^{xxxx}(v_x,0) = \frac{4}{3}\,\frac{T_{\rm rel}\,\eta}{\gamma^4\,(1-c_s^2 v_x^2)^2}
\stp
which agrees with the main-text result for $\kappa^{xxxx}$.

\noindent{\bf Symmetries:}
\begin{align}
\kappa^{ijlm} &= \kappa^{jilm} = \kappa^{ijml} = \kappa^{lmij}
\end{align}
The tensor has only 6 independent components for 2D indices:
$\kappa^{xxxx}$, $\kappa^{xxxy}$, $\kappa^{xxyy}$, $\kappa^{xyxy}$, $\kappa^{xyyy}$, $\kappa^{yyyy}$.

% ─────────────────────────────────────────────────────────────────────────────
\appendix

\section{Implementation notes for the Julia SPH code}
\label{app:impl}

\subsection*{Field and particle notation}

\begin{itemize}
  \item $s_a$ — entropy of particle $a$ (extensive, units: entropy)
  \item $m_a^i$ — momentum of particle $a$ (extensive, units: momentum)
  \item $S(x) = \sum_a s_a\,W(x,x_a)$ — entropy density field (units: entropy/area,
        since $W$ has units $1/\text{area}$ and $\int d^2x\,W = 1$)
  \item $M^i(x) = \sum_a m_a^i\,W(x,x_a)$ — momentum density field; satisfies
        $\int d^2x\,M^i(x) = \sum_a m_a^i$
  \item $\texttt{vol} = \Delta x^2$ — particle volume (area per particle on the
        initialization lattice); appears \emph{only} in initialization to convert
        a density to a per-particle quantity, e.g.\ $s_a = s_{\rm density}\times\texttt{vol}$.
        It does \emph{not} appear in any smoothing sum.
  \item $s_{\rm rel}(e) \equiv (e+p)/T_{\rm rel}$ — relativistic entropy density
        (a thermodynamic field, not a per-particle quantity).
        \emph{Note:} $s_{\rm rel}(e) \neq s_a$.  The approximate relation between
        them is $s_a \approx s_{\rm rel}(e)\,\gamma\,\texttt{vol}$,
        since the per-particle entropy is the entropy density times the
        Lorentz-contracted volume element.
        The specific shear viscosity $\eta/s$ appearing in the transport
        coefficients is properly $\eta/s_{\rm rel}$; in Julia this is
        \texttt{eos.\(\eta\_{\rm by\_s}\)} and the physical viscosity is
        $\eta = (\eta/s_{\rm rel})\,s_{\rm rel}(e)$.
\end{itemize}

\subsection*{Key definitions used in the dissipative step}

\paragraph{The vector $\bbeta$ (\texttt{b\_vec} in Julia).}
The dissipative function $K$ and the pair force $F^i_{\rm visc}(x_a,x_b)$ involve
the 2-vector $\bbeta_a \equiv \dot x_a / \Teff_a$, where
\begin{align}
    \dot x_{i,a} &= \sum_b s_b\,(v/S)_i(x_b)\,W(x_a, x_b)
    \quad\text{(kernel-smoothed velocity)}, \label{eq:dotx}\\
    \Teff_a      &= \E(x_a) - \sum_b s_b\,({\F}/S)(x_b)\,W(x_a, x_b) \,. \label{eq:Teff}
\end{align}
Here $s_b$ is the per-particle entropy (not a density); the sum $\sum_b s_b W$
reproduces the entropy density $S(x_a)$ because $s_b = s_{\rm density}\,\texttt{vol}$
was set at initialization.
Concretely,
\begin{equation}
    \texttt{b\_vec}[a] = \dot x_a \;/\; \Teff_a \,.
\end{equation}

\paragraph{Hamiltonian structure behind $\dot x_a$ and $\Teff_a$.}
Both quantities are Hamilton's equations for the discrete Hamiltonian
\begin{equation}
    {\mathbb H}(s_a, x_a, m_a) = \sum_a s_a\,\E\!\left(M(x_a),\,S(x_a)\right),
\end{equation}
where $M(x_a) = \texttt{pa.M\_smooth}$ and $S(x_a) = \texttt{pa.S\_smooth}$ are the
kernel-smoothed momentum and entropy densities.  The energy per entropy $\E$ is evaluated
as a function of the \emph{smoothed} fields by \texttt{solve\_eos!}:
\begin{equation}
    \E(M_{\rm smooth},\, S_{\rm smooth}) \;=\; \texttt{pa.H\_over\_S}
    \;=\; \frac{(e+p)\gamma^2 - p}{S_{\rm smooth}}\,.
\end{equation}
Its partial derivatives give the thermodynamic conjugates that \texttt{solve\_eos!} stores:
\begin{equation}
    \frac{\partial \E}{\partial M^i_{\rm smooth}} = \frac{v^i}{S}
    = \texttt{pa.v\_over\_S}[i]\,,
    \qquad
    \frac{\partial \E}{\partial S_{\rm smooth}} = -\F/S = -\texttt{pa.F\_over\_S}\,.
\end{equation}
Hamilton's equations then reproduce eqs.~\eqref{eq:dotx}--\eqref{eq:Teff} exactly:
\begin{align}
    \dot x_{i,a} &= \frac{\partial {\mathbb H}}{\partial m_a^i}
    = \sum_b s_b\,\texttt{pa\_b.v\_over\_S}[i]\,W(x_a,x_b), \\
    \Teff_a &= \frac{\partial {\mathbb H}}{\partial s_a}
    = \texttt{pa.H\_over\_S} - \sum_b s_b\,\texttt{pa\_b.F\_over\_S}\,W(x_a,x_b)\,.
\end{align}
\textbf{Key point:} \texttt{pa.H\_over\_S} is \emph{not} the energy per entropy at the
bare particle position; it is $\E$ evaluated at the \emph{smoothed} state
$(M_{\rm smooth}, S_{\rm smooth})$.  Likewise, \texttt{pa.v\_over\_S} is
$\partial\E/\partial M_{\rm smooth}$, computed inside \texttt{solve\_eos!} via the
thermodynamic identity $d\E = (v/S)\,dM - (\F/S)\,dS$.

\paragraph{Gotcha~1: $\bbeta \neq v/T$.}
In relativistic kinetic theory one defines a thermodynamic inverse-temperature
4-vector $\beta^\mu = u^\mu/T_{\rm rel}$.  Its spatial part is
\begin{equation}
    \beta_i = \frac{u_i}{T_{\rm rel}} = \frac{\gamma v_i}{T_{\rm rel}}
            = \frac{v_i}{T_{\rm rel}/\gamma} = \frac{v_i}{\texttt{pa.T}}\,.
\end{equation}
This is \emph{not} the same as $\bbeta_i$ above.  The two coincide only when
(i) the SPH kernel average equals the pointwise value, i.e.\ the field is smooth
on the scale $h$, \emph{and} (ii) $\Teff = T$.  In general $\bbeta$ contains
kernel-smoothing corrections through $\dot x$ (eq.~\eqref{eq:dotx}) and through
$\Teff$ (eq.~\eqref{eq:Teff}).  Using $v_i/\texttt{pa.T}$ in place of $\bbeta_i$
in the pair force would violate energy conservation at $O(h^2)$.

\paragraph{Gotcha~2: $\dot x$ vs.\ the EOS velocity.}
After \texttt{smooth\_to\_particles!} and \texttt{update\_thermodynamics!},
the particle holds $\texttt{pa.v\_over\_S} = v_i/S$ (the local
thermodynamic 3-velocity divided by the local entropy density).
The SPH-smoothed velocity $\dot x_a$ is \emph{not} simply
$\texttt{pa.v\_over\_S}[i] \times \texttt{pa.S\_smooth}$; it is a kernel
average over the neighbors' velocities (eq.~\eqref{eq:dotx}), exactly as computed
in \texttt{hamiltonian\_velocity!}.  In the function \texttt{compute\_b\_vec!}
in \texttt{celllist.jl} both sums (for $\dot x$ and for
$\sum_b s_b\,(F/S)(x_b)\,W$) are evaluated in the same neighbor loop.

\paragraph{Gotcha~3: $T_{\rm rel}\kappa^{ijlm}$ uses the local thermodynamic temperature.}
In the pair force
\begin{equation}
    F^i_{\rm visc}(x_a,x_b) =
    \tfrac{1}{4}s_a s_b\,D_j(a,b)
    \left[\frac{T_{\rm rel}\kappa^{ijlm}}{S^2}\bigg|_a +
          \frac{T_{\rm rel}\kappa^{ijlm}}{S^2}\bigg|_b\right]
    D_l(a,b)\,(\bbeta_{m,b}-\bbeta_{m,a}),
\end{equation}
the factor $T_{\rm rel}\kappa^{ijlm}$ is returned by \texttt{kappa\_tensor(pa, eos, Kappa3D())}
and uses the \emph{local} thermodynamic temperature
$T_{\rm rel} = \texttt{pa.T}\times\texttt{pa.}\gamma$.  This is correct and separate
from the $\Teff$ appearing in $\bbeta$.

\section{Sound waves test}

For a constant background and small fluctuations we can neglect the differences between
$E = (e + p) \gamma^2 - p \simeq e$ and $s_{\rm rel} \simeq S$, so that
$\delta e(t, x) \simeq \frac{(e + p)}{S}\,\delta S \simeq T\,\delta S(t, x)$.
Substituting $\delta e(t, x) = A \cos(kx - \omega t)$ and $M^x(t, x) = A c_s \cos(kx - \omega t)$
into the linearized conservation laws,
\begin{align}
  \partial_t \delta e(t,x) +  \partial_x M^x(t, x) =& 0 \\
  \partial_t M^x(t, x) +  c_s^2\, \partial_x \delta e(t, x) =& 0
\end{align}
The solution has $\omega = c_s k$. 

When dissipation is included the linearized equations become
\begin{align}
  \partial_t \delta e +  \partial_x M^x =& \;0 \,, \\
  \partial_t M^x +  c_s^2\, \partial_x \delta e - \Gamma_s\, \partial_x^2 M^x =& \;0 \,,
\end{align}
where the sound-attenuation coefficient is
\begin{equation}
  \Gamma_s \;=\; \frac{\tfrac{4}{3}\eta + \zeta}{e_0 + p_0}
            \;=\; \frac{\tfrac{4}{3}(\eta/s) + (\zeta/s)}{T_{{\rm rel},0}} \,.
\end{equation}
Inserting the plane-wave ansatz $\delta e, M^x \propto e^{i(kx-\omega t)}$
gives the exact dispersion relation
\begin{equation}
  \omega^2 + i\,\Gamma_s k^2\,\omega - c_s^2 k^2 = 0 \,.
\end{equation}
Solving exactly,
\begin{equation}
\label{eq:sound_dispersion}
  \omega \;=\; \pm\,\omega_R \;-\; i\,\frac{\Gamma_s k^2}{2} \,, \qquad
  \omega_R \;\equiv\; \frac{k}{2}\sqrt{4c_s^2 - \Gamma_s^2 k^2}
           \;=\; c_s k\sqrt{1 - \frac{\Gamma_s^2 k^2}{4c_s^2}} \,,
\end{equation}
valid in the underdamped regime $\Gamma_s k < 2c_s$.
Both modes share the same decay rate $\Gamma_s k^2/2$; the phase velocity
$\omega_R/k$ is shifted below $c_s$ by a term of order ${(\Gamma_s k/c_s)}^2$.
The decay timescale is $\tau_s = 2/(\Gamma_s k^2)$ regardless of whether the
approximation $\Gamma_s k \ll c_s$ holds.

The right-moving eigenmode initialized at $t=0$ with $\delta e(0,x)=A\cos(kx)$
evolves as
\begin{align}
\label{eq:sound_sol}
  \delta e(t,x) &= A\cos(kx - \omega_R t)\,e^{-\Gamma_s k^2 t/2} \,, \\
  M^x(t,x)      &= A\,e^{-\Gamma_s k^2 t/2}
    \left[\frac{\omega_R}{k}\cos(kx-\omega_R t)
         +\frac{\Gamma_s k}{2}\sin(kx-\omega_R t)\right] \,.
\end{align}
The second line follows from the continuity equation
$\partial_t\delta e + \partial_x M^x = 0$, which gives
$M^x_0 = \mathrm{Re}[(\omega/k)\,\delta e_0]$ with the \emph{complex}
$\omega/k = \omega_R/k - i\,\Gamma_s k/2$.
A useful identity: $|\omega/k|^2 = (\omega_R/k)^2 + (\Gamma_s k/2)^2 = c_s^2$
(from the dispersion relation~\eqref{eq:sound_dispersion}), so
\begin{equation}
  M^x(t,x) = A\,c_s\,e^{-\Gamma_s k^2 t/2}\cos(kx-\omega_R t - \varphi)\,,
  \qquad
  \varphi = \arctan\!\left(\frac{\Gamma_s k^2}{2\,\omega_R}\right) \,.
\end{equation}
The phase lag $\varphi$ vanishes as $\Gamma_s\to 0$.

\paragraph{Initial condition.}
Using $\delta e \simeq T_{{\rm rel},0}\,\delta S$ the IC for the right-moving
eigenmode is
\begin{align}
\label{eq:sound_ic}
  \delta S(x) &= \frac{A}{T_{{\rm rel},0}}\cos(kx) \,, \\
  M^x(x)      &= A\left[\frac{\omega_R}{k}\cos(kx)
                        +\frac{\Gamma_s k}{2}\sin(kx)\right]
               = A\,c_s\cos(kx - \varphi) \,.
\end{align}
For $\eta/s = \zeta/s = 0$ one has $\Gamma_s = 0$, $\omega_R = c_s k$,
$\varphi = 0$, and the IC reduces to the one in \texttt{testsound.jl}.
For the viscous test one sets $\eta/s > 0$ and uses the GENERIC integrator
\texttt{generic\_step!}.

\section{ Dissipative force per volume}

We defined
\begin{align}
  M(x) = \sum m_a W(x, x_a)
\end{align}
The force per volume is  
\begin{align}
     \frac{\partial M(x)}{\partial t} =  \sum_a  \frac{dm_a}{dt} W(x, x_a)   -  m_a \dot x_a  \partial_x W(x, x_a)  
\end{align}
For a small amplitude shear wave in a uniform background
\begin{align}
  M^y (x) = A \cos(kx) 
\end{align}
We have  the following estimate:
\begin{align}
  \partial_t M^y(x) = - \partial_x \Pi^{yx}  \simeq  \partial_x  (T_{\rm rel }\eta \, \partial_x \beta^y) \simeq  \frac{\eta}{e + p } \partial_x^2 M^y = \frac{\eta/s}{T_0} \partial_x^2 M^y
\end{align}


% ─────────────────────────────────────────────────────────────────────────────
\section{Two-level dissipative loop for the Julia SPH code}
\label{app:dissipative_loop}

\subsection*{The problem with the original pair-based discretization}

The dissipative function originally proposed in the main text is
\begin{equation}
    K = \frac{1}{16}\sum_{ab} s_a s_b\,
    D_j(x_a,x_b)(\bbeta_{i,b}-\bbeta_{i,a})
    \left[\frac{T_{\rm rel}\kappa^{ijlm}}{S^2}\bigg|_a
         +\frac{T_{\rm rel}\kappa^{ijlm}}{S^2}\bigg|_b\right]
    D_l(x_a,x_b)(\bbeta_{m,b}-\bbeta_{m,a})\,.
\end{equation}
In the continuum limit this should approximate
\begin{equation}
    K \;\longrightarrow\; \frac{1}{2}\int d^2x\; \partial_j\beta_i\,
    T_{\rm rel}\kappa^{ijlm}\,\partial_l\beta_m \,.
\end{equation}
The inconsistency is that $D_j(x_a,x_b)$ and $D_l(x_a,x_b)$ are both kernel
gradients evaluated at the \emph{same pair} $(a,b)$.  Their supports overlap
completely, so the pair sum does not factorize into two independent gradient
estimates at a point.  In the continuum limit the product of two coincident
kernel gradients produces a $\delta(0)$ divergence, and the normalization is
wrong.  Moreover, the pair contraction $D_j(a,b)D_l(a,b)\Delta\bbeta_j\Delta\bbeta_l$
captures only the \emph{longitudinal} component of the strain along $\hat r_{ab}$;
transverse components of $\partial_j\beta_i$ are missed at $\mathcal{O}(h^2)$.

\subsection*{The two-level remedy}

The fix is to compute the strain tensor $\varepsilon_{ji}$ at each particle site
\emph{first}, using a proper SPH gradient estimator, and then to build the force
from the stored strains.  This separates the two kernel gradients into independent
neighbor sums, eliminating the overlap.

\paragraph{Discretize the continuum dissipative function.}
Define the SPH strain at particle $a$ as
\begin{equation}
\label{eq:strain-impl}
    \varepsilon_{ji}(x_a) \;\equiv\;
    \frac{1}{S(x_a)}\sum_b s_b\,
    \bigl(\bbeta_{i,b} - \bbeta_{i,a}\bigr)\,D_j(x_a,x_b) \,,
\end{equation}
which is an antisymmetric SPH gradient estimator for $\partial_j\beta_i$ at
particle $a$.  The discrete dissipative function is then
\begin{equation}
\label{eq:K_twolevel-impl}
    K \;=\; \frac{1}{2} \sum_a \frac{s_a}{S(x_a)}\;
    \varepsilon_{ji}(x_a)\,
    T_{\rm rel}\kappa^{ijlm}\big|_a\;
    \varepsilon_{lm}(x_a) \,.
\end{equation}
Here the two factors of $D_j$ live in \emph{separate} neighbor sums (one for
each index of $\varepsilon$), so there is no overlap.  The volume element
$s_a/S(x_a)$ is the correct SPH approximation to $d^2x$.

\paragraph{Viscous force from $\partial K/\partial \bbeta_{i,c}$.}
Taking the derivative of $K$ with respect to $\bbeta_{i,c}$ and using
$f^i_{\rm visc}(x_c) = -\partial K/\partial\bbeta_{i,c}$ yields two terms,
depending on whether $c$ appears as the reference particle (the $a=c$
term) or as a neighbor (the $b=c$ term in the sum over $b$):

\begin{equation}
\label{eq:force_twolevel}
    f^i_{\rm visc}(x_c) =
    \underbrace{+\frac{s_c}{S_c^2}\sum_b s_b\,D_j(c,b)\,
    T_{\rm rel}\kappa^{ijlm}\big|_c\;\varepsilon_{lm}(x_c)}_{\text{term 1: }c\text{ is reference particle}}
    \;-\;
    \underbrace{\sum_a \frac{s_a s_c}{S_a^2}\,D_j(a,c)\,
    T_{\rm rel}\kappa^{ijlm}\big|_a\;\varepsilon_{lm}(x_a)}_{\text{term 2: }c\text{ is neighbor of }a}\,.
\end{equation}
This can be written
\begin{equation}
\label{eq:force_combined}
    f^i_{\rm visc}(x_c) =
     \sum_a s_c s_a \,D_j(c,a)\,
     \left( \left. \frac{T_{\rm rel}\kappa^{ijlm}}{S^2}\right|_c\;\varepsilon_{lm}(x_c)
    \;+\;
    \left. \frac{T_{\rm rel}\kappa^{ijlm}}{S^2}\right|_a\;\varepsilon_{lm}(x_a) \right)
\end{equation}
Finally relabelling we find
\begin{equation}
\label{eq:force_relabeled}
    f^i_{\rm visc}(x_a) =
     \sum_b
      s_a s_b \,D_j(x_a, x_b)\,
     \left( \left. \frac{T_{\rm rel}\kappa^{ijlm}}{S^2}\right|_a\;\varepsilon_{lm}(x_a)
    \;+\;
    \left. \frac{T_{\rm rel}\kappa^{ijlm}}{S^2}\right|_b\;\varepsilon_{lm}(x_b) \right)
\end{equation}

Momentum conservation $\sum_c f^i_{\rm visc}(x_c) = 0$ follows from the
antisymmetry of $D_j(x_a, x_b)$ . We write this as in the main text
\begin{equation}
   \left(\frac{\Pi^{ij}}{S^2}(x_a)  
   +
   \frac{\Pi^{ij}}{S^2}(x_b)  
   \right)
     =
     -\left( \left. \frac{T_{\rm rel}\kappa^{ijlm}}{S^2}\right|_a\;\varepsilon_{lm}(x_a)
    \;+\;
    \left. \frac{T_{\rm rel}\kappa^{ijlm}}{S^2}\right|_b\;\varepsilon_{lm}(x_b) \right)
\end{equation}
so
\begin{equation}
\label{eq:force_Pi}
    f^i_{\rm visc}(x_a) =
     -\sum_b s_a s_b \,D_j(x_a, x_b)\,
   \left(\frac{\Pi^{ij}}{S^2}(x_a)  
   +
   \frac{\Pi^{ij}}{S^2}(x_b)  
   \right)
\end{equation}


Entropy production $\dot S = 2K \geq 0$ is guaranteed because $K$ is a
positive-definite quadratic form in $\varepsilon$ and $\kappa^{ijlm}$ is
positive definite.

\paragraph{Entropy update.}
Once $f^i_{\rm visc}(x_c)$ is known, the entropy is updated from energy
conservation (eq.~50 of the main text):
\begin{equation}
\label{eq:entropy_update}
    \Delta s_a = -\bbeta_a \cdot \Delta t\, f_{\rm visc,a} \,.
\end{equation}

\subsection*{Implementation in Julia — three passes}

The two-level force requires two neighbor loops plus one particle loop.
The total cost is approximately $3\times$ the standard pair-force loop,
but the continuum limit is correct.

\paragraph{Pass 1 — compute strain tensor at every particle.}
Loop over ALL neighbors $b \neq a$ (not just $b > a$, since
$\varepsilon_{ji}(x_a)$ needs contributions from every neighbor):
\begin{lstlisting}
for a in 1:N
    ε[a] = zero(SMatrix{2,2,Float64})
    for b in neighbors(a)            # all b != a
        r_mic, r_mag = minimum_image(x[a]-x[b], L)
        r_mag >= 2h && continue
        D_ab  = wendland_c2_grad(r_mic, r_mag, h)
        Δb    = b_vec[b] - b_vec[a]
        ε[a] += s[b] * (Δb * D_ab^T) / S_smooth[a]   # outer product
    end
end
\end{lstlisting}

\paragraph{Pass 2 — compute viscous force from stored strains.}
Loop over pairs $b > a$ and apply Newton's third law:
\begin{lstlisting}
for a in 1:N, b in neighbors(a) with b > a
    r_mic, r_mag = minimum_image(x[a]-x[b], L)
    r_mag >= 2h && continue
    D_ab   = wendland_c2_grad(r_mic, r_mag, h)
    κ_a    = kappa_tensor(particles[a], eos)     # T_rel κ^{ijlm} at a
    κ_b    = kappa_tensor(particles[b], eos)
    w_a    = s[a] * s[b] / S_smooth[a]^2
    w_b    = s[a] * s[b] / S_smooth[b]^2
    f_ab   = w_a * κ_a * ε[a] * D_ab
           - w_b * κ_b * ε[b] * D_ab
    f_visc[a] += f_ab
    f_visc[b] -= f_ab                            # Newton's third law
end
\end{lstlisting}

\paragraph{Pass 3 — entropy update.}
One loop over particles, no neighbor search required:
\begin{lstlisting}
for a in 1:N
    Δs[a] = -dot(b_vec[a], dt * f_visc[a])
    s[a] += Δs[a]
end
\end{lstlisting}

\subsection*{Checks to run after implementation}

\begin{enumerate}
    \item \textbf{Momentum conservation:} $\sum_a f^i_{\rm visc}(x_a) = 0$
          to machine precision after Pass~2.
    \item \textbf{Second law:} $\sum_a \Delta s_a \geq 0$ after Pass~3.
    \item \textbf{Energy conservation:} $|\sum_a \Teff_a\,\Delta s_a
          + \dot x_a \cdot \Delta t\, f_{\rm visc,a}| < \epsilon_{\rm tol}$
          (GENERIC degeneracy condition).
    \item \textbf{Shear wave damping:} with $\eta > 0$, $\zeta = 0$,
          a transverse wave $M^y(x) = A\cos(kx)$ should damp at rate
          $\gamma_k = (\eta/s)\,k^2/T_0$ (Appendix~D of main text).
\end{enumerate}


\section{Generating noise}

We first generate a \emph{unit-normalized} noise $\hat\xi^{ij}$ whose correlation
function equals the bare kappa tensor,
\begin{align}
\label{eq:xihat_cov}
 \llangle \hat\xi^{ij} \hat\xi^{lm} \rrangle = \kappa^{ijlm} \, ,
\end{align}
with $\kappa^{ijlm}$ the tensor built from $2\eta(\dots)+\zeta(\dots)$ below
(no factor of $T_{\rm rel}$).  The physical density-frame noise is then the
rescaling
\begin{align}
\label{eq:xi_from_xihat}
 \boxed{\;\xi^{ij} \;=\; \sqrt{2\,T_{\rm rel}}\;\hat\xi^{ij}\;}
\end{align}
which satisfies the usual fluctuation--dissipation relation
$\llangle \xi^{ij}(x,t)\,\xi^{lm}(x',t')\rrangle = 2\,T_{\rm rel}\,\kappa^{ijlm}(v)\,
\delta(t-t')\,\delta^d(x-x')$, matching $\Pi^{ij} = -T_{\rm rel}\kappa^{ijlm}\partial_l\beta_m$.
The construction of $\hat\xi^{ij}$ is done for every particle as follows:
\begin{enumerate}
\item Produce the Landau noise in the local rest frame  
$\xi^{\mu\nu}_{LRF} $ (this is a symmetric tensor purely spatial).  Any component of this tensor 
with a zero index is zero
\begin{align}
 \xi^{0i}_{LRF} = \xi^{i0}_{LRF} = \xi^{00} = 0
\end{align}

We need the covariance matrix of $\xi^{ij}_{LRF}$ to satisfy 
\begin{equation}
  \llangle \xi^{ij}_{LRF} \xi^{lm}_{LRF} \rrangle = 2 \eta \left[ \delta^{(il} \delta^{j)m}  -\frac{1}{d} \delta^{ij} \delta^{lm} \right]  + \zeta  \delta^{ij} \delta^{lm} 
\end{equation}
One way to do this is to pick one tensor
\begin{align}
\xi_{L}^{ij} = \sqrt{d (d + 2) \eta} \,  \xi_1 \left(n^i n^j - \tfrac{1}{d} \delta^{ij} \right)   +  \sqrt{\zeta} \, \xi_2 \delta^{ij} 
\end{align}
Here $\xi_1$ and $\xi_2$ are independent random numbers with variance one,
and $n^i$ is a unit vector drawn uniformly from the sphere ($d=3$) or circle ($d=2$).
The relevant moments of the uniform measure are
\begin{align}
  \llangle n^i n^j \rrangle &= \frac{1}{d}\,\delta^{ij} \\
  \llangle n^i n^j n^l n^m \rrangle &= \frac{1}{d(d+2)}\left(\delta^{ij}\delta^{lm} + \delta^{il}\delta^{jm} + \delta^{im}\delta^{jl}\right)
\end{align}
One can then verify directly that the covariance is correct.  Using the cross terms
\begin{align}
\llangle \left(n^i n^j - \tfrac{1}{d}\delta^{ij}\right)\!\left(n^l n^m - \tfrac{1}{d}\delta^{lm}\right)\rrangle
&= \frac{\delta^{il}\delta^{jm}+\delta^{im}\delta^{jl}}{d(d+2)}
  - \frac{2}{d^2(d+2)}\,\delta^{ij}\delta^{lm}
\end{align}
one finds
\begin{align}
\llangle \xi^{ij}_L \xi^{lm}_L \rrangle
&= d(d+2)\eta\left[\frac{\delta^{il}\delta^{jm}+\delta^{im}\delta^{jl}}{d(d+2)}
   - \frac{2\,\delta^{ij}\delta^{lm}}{d^2(d+2)}\right]
  + \zeta\,\delta^{ij}\delta^{lm} \nonumber\\
&= 2\eta\left[\delta^{(il}\delta^{j)m} - \frac{1}{d}\delta^{ij}\delta^{lm}\right]
  + \zeta\,\delta^{ij}\delta^{lm}
\end{align}
as required.


\item Boost to the lab frame.
The fluid cell has four-velocity $u^\mu = (\gamma, \gamma v^i)$ in the lab frame.
The Lorentz boost that takes the LRF rest vector $(1,0,0,0)$ to $u^\mu$
has matrix elements
\begin{align}
  \Lambda^0{}_0 &= \gamma, &
  \Lambda^0{}_j &= \gamma\, v_j, \\
  \Lambda^i{}_0 &= \gamma\, v^i, &
  \Lambda^i{}_j &= \delta^i_j + \frac{\gamma - 1}{v^2}\,v^i v_j,
\end{align}
where $v^2 = v_i v^i$.
The Landau-frame noise tensor in the lab frame is then
\begin{align}
 \xi^{\mu\nu}_{L} = \Lambda^{\mu}{}_{\alpha}\, \Lambda^{\nu}{}_{\beta}\, \xi^{\alpha\beta}_{LRF}\,.
\end{align}
Since $\Delta^{\mu\nu} = g^{\mu\nu} + u^\mu u^\nu$ (with $g^{\mu\nu} = \mathrm{diag}(-1,1,1,1)$)
is the projector onto the space orthogonal to $u^\mu$, one can verify that
\begin{align}
  \llangle \xi^{\mu\nu}_{L}\, \xi^{\alpha\beta}_{L}  \rrangle
  = 2\eta\left(\Delta^{(\mu\alpha}\Delta^{\nu)\beta}
    - \frac{1}{d}\Delta^{\mu\nu}\Delta^{\alpha\beta}\right)
    + \zeta\,\Delta^{\mu\nu}\Delta^{\alpha\beta}\,.
\end{align}


\item Apply the frame-change matrix $\kappa_{\mu\nu}^{ij}$ to obtain the
unit-normalized Density-Frame noise:
\begin{align}
   \hat\xi^{ij}  = \kappa_{\mu\nu}^{ij}\, \xi^{\mu\nu}_{L}\,.
\end{align}
The matrix $\kappa^{ij}_{\mu\nu}$ is built from the comoving-frame
basis vectors and the equation of state.  Define
\begin{align}
  e^i_\mu &= \delta^i_\mu - v^i \delta^0_\mu \,, \label{eq:e_def}\\
  h^{ij}  &= e^i_\mu e^j_\nu \eta^{\mu\nu} = \delta^{ij} - v^i v^j \,. \label{eq:h_def}
\end{align}
Here $e^i_\mu$ projects a lab-frame covariant index onto the comoving spatial
directions, and $h^{ij}$ is the induced spatial metric in the comoving frame.
The transformation matrix is (theory.pdf eq.\ (43))
\begin{align}
  \kappa^{ij}_{\mu\nu}
  = e^i_{(\mu} e^j_{\nu)}
    - h^{ij}\,\frac{\partial\mathcal{P}}{\partial\mathcal{T}^{0\rho}}
      \delta^0_{(\mu}\delta^\rho_{\nu)} \,, \label{eq:kappa_def}
\end{align}
where the pressure derivative (theory.pdf eq.\ (44)) is
\begin{align}
  \frac{\partial\mathcal{P}}{\partial\mathcal{T}^{0\rho}}\,
  \delta^0_\mu\delta^\rho_\nu
  = -\frac{c_s^2}{1 - v^2 c_s^2}
    \left[\frac{1}{\gamma^2}\Delta^0_\mu\Delta^0_\nu - u_\mu u_\nu\right] \,.
    \label{eq:dPdT}
\end{align}
Here $\Delta^{\mu\nu} = g^{\mu\nu} + u^\mu u^\nu$ is the projector orthogonal to
$u^\mu$, with $\Delta^0_\mu \equiv g^{0\nu}\Delta_{\nu\mu}$.
It is also useful to introduce the scalar projector (theory.pdf eq.\ (45))
\begin{align}
  P_{\mu\nu} = -c_s^2\,u_\mu u_\nu + \frac{1}{d}\,\Delta_{\mu\nu} \,,
  \label{eq:Pmunu}
\end{align}
in terms of which $\kappa^{ij}_{\mu\nu}$ decomposes into shear and bulk parts
(theory.pdf eq.\ (46)):
\begin{align}
  \kappa^{ij}_{\mu\nu}
  &= \underbrace{\left(e^i_{(\mu}e^j_{\nu)}
      - \frac{h^{ij}}{\langle Ph\rangle}(ePe)_{\mu\nu}\right)}_{\mathring\kappa^{ij}_{\mu\nu}}
    + \underbrace{\frac{h^{ij}}{\langle Ph\rangle}\,P_{\mu\nu}}_{\bar\kappa^{ij}_{\mu\nu}} \,,
\end{align}
where
\begin{align}
  (ePe)_{\mu\nu} &= e^i_\mu P_{ij} e^j_\nu
    = -\frac{c_s^2}{\gamma^2}\,\Delta^0_\mu\Delta^0_\nu
      + \frac{1}{d}\,\Delta_{\mu\nu} \,, \\
  \langle Ph\rangle &= P_{ij}\,h^{ji} = 1 - c_s^2 v^2 \,.
\end{align}
The covariance of the resulting unit-normalized Density-Frame noise then
satisfies $\llangle\hat\xi^{ij}\hat\xi^{lm}\rrangle = \kappa^{ijlm}$ with
(theory.pdf eq.\ (54))
\begin{align}
  \kappa^{ijlm} = \kappa^{ij}_{\mu\nu}
    \left[2\eta\!\left(\Delta^{(\mu\rho}\Delta^{\nu)\sigma}
      - \tfrac{1}{d}\Delta^{\mu\nu}\Delta^{\rho\sigma}\right)
      + \zeta\,\Delta^{\mu\nu}\Delta^{\rho\sigma}\right]
    \kappa^{lm}_{\rho\sigma} \,,
\end{align}
and the physical noise $\xi^{ij} = \sqrt{2T_{\rm rel}}\,\hat\xi^{ij}$ of
\Eq{eq:xi_from_xihat} then obeys the FDT normalization
$\llangle\xi^{ij}\xi^{lm}\rrangle = 2 T_{\rm rel}\kappa^{ijlm}$.

\end{enumerate}

\section{First stochastic steps}

\subsection{Physical Parameters}
We need to decide on the parameters
\begin{itemize}
\item Shear Length
\begin{align}
 \ell  \equiv  \frac{\eta }{(e/c^2 +  p) c_s}
\end{align}
we note that $sT = e + c^2 p$. 
\item Interparticle spacing 
\begin{align}
 \ell_Q  \equiv  \frac{1}{(s_0)^{1/2}}
\end{align}
\item  Speed of light $c$

\item  Temperature $T_0$
\end{itemize}

We can choose units where $c = T_0 = \ell_Q = 1$.  

\subsection{Defined and Derived Physical Parameters}
Then we derive and define:
\begin{itemize}
\item We define the thermal momentum
\begin{align}
  p_0 \equiv  \frac{T_0}{c}
\end{align}

\item The dimensionless shear viscosity
is the mean free path in these units and reads
\begin{align}
   \hat \eta \equiv \ell /\ell_Q   > 1
\end{align} 

  
\item  The planks constant comes up in the definition of the entropy
density 
\begin{equation}
s_0  \equiv   \frac{ \kappa^2 p_0^2 }{\hbar^2}  \qquad    
\end{equation}
and thus
\begin{align}
  \ell_Q =   \frac{\hbar}{\kappa\, p_0}
\end{align}
For a classical relativistic ideal gas with  $N^2$ degrees of freedom
 in 2d we have
\begin{align}
    \kappa^2 \equiv \frac{N^2}{\pi} 
\end{align}
\end{itemize}

We find (using that $s_0 T_0 = e + c^2 p$) we have
\begin{equation}
 \hat \eta = \left(\frac{\eta}{s \hbar } \right)   \kappa \left(\frac{c}{c_s} \right)
\end{equation}
Numerically 
\begin{equation}
 \kappa =  3.4    \left( \frac{N^2}{37}  \right)^{1/2}
\end{equation}
and  for $c_s = c/2$
\begin{equation}
 \hat \eta = 0.54 \left( \frac{4\pi \eta/s}{\hbar }\right)   \left( \frac{\kappa }{3.4}  \right) 
\end{equation}

\subsection{Basic description}

We are applying hydro over a length $L$. There are viscous corrections 
\begin{align}
      {\rm classical} \sim \frac{\ell}{L}
\end{align}
There are also $1/\sqrt{N}$ corrections 
\begin{align}
    {\rm noise} \sim \frac{1}{\sqrt{L^2 s} } =  \frac{\ell_Q}{L}
\end{align}
which in two dimensions scale similarly to the classical terms. 

The ratio of these corrections is 
\begin{align}
     \frac{ {\rm classical}}{{\rm noise}  } \sim \hat \eta
\end{align}

\subsection{Numerical Parameters}

\begin{itemize}
\item $h$  coarsening length
\item $a$  inter-sph-particle spacing
\end{itemize}

We want to choose $h$ to be of order $\ell$ and $a$ to be of order $\ell_Q$ or 
smaller.  So we choose
\begin{align}
  a\sim \ell_Q   \lesssim    h \sim \ell  \lesssim  L
\end{align}

To get started lets go crazy taking $\hat \eta \sim 10$

\subsection{\texorpdfstring{Relation to the code's $T_{\rm scale}$ normalization}{Relation to the code's Tscale normalization}}

In \texttt{sph.jl}, \texttt{ConformalEOS} parameterizes the temperature via
\begin{align}
  T_{\rm rel}(e) = T_{\rm scale}\, e^{c_s^2/(1+c_s^2)} \,.
\end{align}
For the $d=2$ fluid ($c_s^2=1/2$) this gives $e = (T_{\rm rel}/T_{\rm scale})^3$, and with $p=c_s^2 e$,
\begin{align}
  s = \frac{e+p}{T_{\rm rel}} = \frac{3}{2}\,\frac{T_{\rm rel}^2}{T_{\rm scale}^3} \,.
\end{align}
Comparing to $s_0 = \kappa^2 p_0^2/\hbar^2 = (N^2/\pi)\,T_0^2$ (natural units $\hbar=c=1$, $T_0=T_{\rm rel}$)
and matching the coefficients of $T^2$,
\begin{align}
  T_{\rm scale} = \left(\frac{3\pi}{2N^2}\right)^{1/3} = \frac{(3\pi/2)^{1/3}}{N^{2/3}} \,,
  \qquad\text{equivalently}\qquad
  N^2 = \frac{3\pi}{2\,T_{\rm scale}^3}\,.
\end{align}
Using $\kappa^2 = N^2/\pi$ (line above) to eliminate $N^2$, this simplifies to
\begin{align}
  T_{\rm scale} = \left(\frac{3}{2\kappa^2}\right)^{1/3} = \frac{(3/2)^{1/3}}{\kappa^{2/3}} \,,
  \qquad\text{equivalently}\qquad
  \kappa^2 = \frac{3}{2\,T_{\rm scale}^3}\,,
\end{align}
i.e.\ $T_{\rm scale}$ scales as $\kappa^{-2/3}$, with the $\pi$ absorbed into $\kappa$ leaving
the clean coefficient $(3/2)^{1/3}\approx1.145$.
As a check, the numerical example above ($N^2=37 \Rightarrow \kappa\approx3.4$) gives
$T_{\rm scale}\approx0.50$ — noticeably different from the arbitrary placeholder
$T_{\rm scale}=1.0$ currently used in \texttt{testeos2d.jl}, which only exercises the
EOS-solver mechanics and was never meant to be physically normalized.

\subsection{Checking equilibration}

A first step is to choose and
\end{document}




